\section{习题}

% 答案框：\boxed 的内容处于数学模式，中文需另行包裹且不能自动换行，
% 故用整行宽的 \fbox + \parbox 排版答案。
\providecommand{\chapfourans}[1]{%
  \par\noindent\fbox{\parbox{\dimexpr\linewidth-2\fboxsep-2\fboxrule\relax}{\textbf{答案：}#1}}\par}

\subsection{本章速览}
\begin{itemize}
    \item 期望是“以概率为权的加权平均”：离散型 $E(X)=\sum_i x_ip_i$，连续型 $E(X)=\int_{-\infty}^{+\infty}xf(x)\,\mathrm{d}x$，两者都要求相应的和或积分绝对收敛。$E(X)$ 只由分布决定，与随机变量怎样取值无关。
    \item 期望的线性性无条件成立：$E(aX+bY+c)=aE(X)+bE(Y)+c$。乘积拆开需要条件：$X,Y$ 相互独立时才有 $E(XY)=E(X)E(Y)$。求 $E[g(X)]$ 不必先求 $g(X)$ 的分布，直接用 $\sum_i g(x_i)p_i$ 或 $\int g(x)f(x)\,\mathrm{d}x$。
    \item 方差 $D(X)=E\bigl[(X-E(X))^2\bigr]=E(X^2)-[E(X)]^2$，恒有 $D(aX+b)=a^2D(X)$；一般情形 $D(X\pm Y)=D(X)+D(Y)\pm2\Cov(X,Y)$，仅当 $X,Y$ 不相关时才可相加。
    \item 协方差 $\Cov(X,Y)=E(XY)-E(X)E(Y)$ 是双线性运算，$\Cov(X,X)=D(X)$；相关系数 $\rho_{XY}=\dfrac{\Cov(X,Y)}{\sqrt{D(X)D(Y)}}$ 满足 $|\rho_{XY}|\le1$，它只度量线性相关程度。
    \item 独立一定不相关，不相关不一定独立；但 $(X,Y)$ 服从二维正态分布时，不相关与独立等价。$|\rho_{XY}|=1$ 当且仅当 $X,Y$ 以概率 $1$ 成线性关系。
    \item $k$ 阶原点矩为 $E(X^k)$，$k$ 阶中心矩为 $E[(X-E(X))^k]$，方差就是二阶中心矩。条件期望 $E(X\mid Y=y)$ 是 $y$ 的函数，满足重期望公式 $E(X)=E[E(X\mid Y)]$ 与方差分解 $D(X)=E[D(X\mid Y)]+D[E(X\mid Y)]$。
\end{itemize}

\subsection{基础过关题}

\textbf{1.} 判断下列命题是否正确，并简述理由。
\begin{enumerate}[label=(\arabic*)]
    \item 若 $E(X^2)=0$，则 $P(X=0)=1$；
    \item 若 $D(X+Y)=D(X)+D(Y)$，则 $X$ 与 $Y$ 相互独立；
    \item 若 $X$ 与 $Y$ 不相关，则 $D(X-Y)=D(X)+D(Y)$；
    \item 若 $(X,Y)$ 服从二维正态分布且 $\rho_{XY}=0$，则 $X$ 与 $Y$ 相互独立；
    \item 若 $E(X)=E(Y)$ 且 $D(X)=D(Y)$，则 $X$ 与 $Y$ 同分布。
\end{enumerate}

\ifshowsolutions%
\boxed{\textbf{解}}
\textbf{思路：}逐条回到定义。涉及“独立”“不相关”的命题，要分清哪个条件强；否命题只需给一个反例。

(1) 正确。$X^2\ge0$，若非负随机变量的期望为零，则它几乎处处为零，即 $P(X^2=0)=1$，也就是 $P(X=0)=1$。

(2) 错误。展开得 $D(X+Y)=D(X)+D(Y)+2\Cov(X,Y)$，题设只说明 $\Cov(X,Y)=0$，即 $X,Y$ 不相关；不相关推不出独立。反例：$X\sim N(0,1)$，$Y=X^2$，则 $\Cov(X,Y)=E(X^3)-E(X)E(X^2)=0$，但 $Y$ 由 $X$ 完全确定，二者不独立。

(3) 正确。$D(X-Y)=D(X)+D(Y)-2\Cov(X,Y)$，不相关即 $\Cov(X,Y)=0$，故 $D(X-Y)=D(X)+D(Y)$。

(4) 正确。二维正态分布由两个边际分布和相关系数唯一确定；$\rho_{XY}=0$ 时联合密度等于两个边际密度之积，因此独立。这是“不相关推独立”的一个重要例外。

(5) 错误。期望与方差只刻画位置和离散程度，不能决定分布形状。反例：$X\sim N(0,1)$，$Y$ 以各 $1/2$ 的概率取 $-1,1$，则 $E(X)=E(Y)=0$，$D(X)=D(Y)=1$，但 $X$ 连续、$Y$ 离散，分布不同。

\chapfourans{(1) 正确\quad(2) 错误\quad(3) 正确\quad(4) 正确\quad(5) 错误}

\else \answerarea[8] \fi

\textbf{2.} 设离散型随机变量 $X$ 的分布律为
\[
\begin{array}{c|cccc}
X & -2 & -1 & 0 & 1 \\
\hline
P & 0.1 & 0.3 & 0.4 & 0.2
\end{array}
\]
\begin{enumerate}[label=(\arabic*)]
    \item 求 $E(X)$、$E(X^2)$ 与 $D(X)$；
    \item 求 $E(|X|)$；
    \item 求 $E(2X+1)$ 与 $D(2X+1)$。
\end{enumerate}

\ifshowsolutions%
\boxed{\textbf{解}}
\textbf{思路：}先按定义算 $E(X)$ 与 $E(X^2)$，再用 $D(X)=E(X^2)-[E(X)]^2$；第 (2) 问用 $E[g(X)]=\sum_i g(x_i)p_i$，不要误以为 $E(|X|)=|E(X)|$。

(1)
\[
E(X)=(-2)(0.1)+(-1)(0.3)+0(0.4)+1(0.2)=-0.2-0.3+0.2=-0.3.
\]
\[
E(X^2)=4(0.1)+1(0.3)+0(0.4)+1(0.2)=0.4+0.3+0.2=0.9.
\]
\[
D(X)=E(X^2)-[E(X)]^2=0.9-(-0.3)^2=0.9-0.09=0.81.
\]

(2) 按 $g(x)=|x|$ 逐项加权：
\[
E(|X|)=2(0.1)+1(0.3)+0(0.4)+1(0.2)=0.2+0.3+0.2=0.7.
\]
注意 $|E(X)|=0.3\ne0.7$，说明取绝对值与取期望一般不能交换。

(3) 由线性性 $E(2X+1)=2E(X)+1=2(-0.3)+1=0.4$；由 $D(aX+b)=a^2D(X)$ 得
\[
D(2X+1)=2^2D(X)=4\times0.81=3.24.
\]

\chapfourans{(1) $E(X)=-0.3$，$E(X^2)=0.9$，$D(X)=0.81$；(2) $E(|X|)=0.7$；(3) $E(2X+1)=0.4$，$D(2X+1)=3.24$}

\else \answerarea[8] \fi

\textbf{3.} 设连续型随机变量 $X$ 的概率密度为
\[
f(x)=\begin{cases}
ax^2, & 0<x<2,\\
0, & \text{其他},
\end{cases}
\]
其中 $a>0$。
\begin{enumerate}[label=(\arabic*)]
    \item 求常数 $a$；
    \item 求 $E(X)$ 与 $D(X)$；
    \item 求 $P\!\left(\frac12<X<\frac32\right)$。
\end{enumerate}

\ifshowsolutions%
\boxed{\textbf{解}}
\textbf{思路：}先用规范性 $\int_{-\infty}^{+\infty}f(x)\,\mathrm{d}x=1$ 定出 $a$；期望与方差按 $\int xf(x)\,\mathrm{d}x$、$\int x^2f(x)\,\mathrm{d}x$ 计算；概率直接对密度积分。

(1) 由规范性
\[
1=\int_0^2 ax^2\,\mathrm{d}x=a\cdot\frac{8}{3},
\]
故 $a=\dfrac{3}{8}$。

(2)
\[
E(X)=\int_0^2 x\cdot\frac38x^2\,\mathrm{d}x=\frac38\int_0^2x^3\,\mathrm{d}x=\frac38\cdot4=\frac32,
\]
\[
E(X^2)=\int_0^2 x^2\cdot\frac38x^2\,\mathrm{d}x=\frac38\int_0^2x^4\,\mathrm{d}x=\frac38\cdot\frac{32}{5}=\frac{12}{5}.
\]
于是
\[
D(X)=\frac{12}{5}-\left(\frac32\right)^2=\frac{12}{5}-\frac94=\frac{48-45}{20}=\frac{3}{20}.
\]

(3)
\[
P\!\left(\frac12<X<\frac32\right)
=\frac38\int_{1/2}^{3/2}x^2\,\mathrm{d}x
=\frac38\cdot\frac{(3/2)^3-(1/2)^3}{3}
=\frac38\cdot\frac{13}{12}
=\frac{13}{32}.
\]

\chapfourans{(1) $a=\dfrac38$；(2) $E(X)=\dfrac32$，$D(X)=\dfrac{3}{20}$；(3) $P\!\left(\dfrac12<X<\dfrac32\right)=\dfrac{13}{32}$}

\else \answerarea[8] \fi

\textbf{4.} 填空题：利用常见分布的期望与方差反求参数。
\begin{enumerate}[label=(\arabic*)]
    \item 设 $X\sim B(n,p)$，若 $E(X)=6$，$D(X)=2.4$，则 $n=\underline{\qquad}$，$p=\underline{\qquad}$；
    \item 设 $X\sim P(\lambda)$，若 $E(X^2)=6$，则 $\lambda=\underline{\qquad}$；
    \item 设 $X\sim U(a,b)$，若 $E(X)=2$，$D(X)=\dfrac43$，则 $a=\underline{\qquad}$，$b=\underline{\qquad}$；
    \item 设 $X\sim E(\lambda)$，若 $E(X^2)=8$，则 $\lambda=\underline{\qquad}$，$D(X)=\underline{\qquad}$。
\end{enumerate}

\ifshowsolutions%
\boxed{\textbf{解}}
\textbf{思路：}写出各分布的 $E(X)$、$D(X)$ 关于参数的表达式，联立求解，并注意参数范围（$0<p<1$，$\lambda>0$，$a<b$）。

(1) $E(X)=np=6$，$D(X)=np(1-p)=2.4$。两式相除得 $1-p=0.4$，故 $p=0.6$，$n=6/0.6=10$。

(2) 泊松分布 $E(X)=D(X)=\lambda$，故 $E(X^2)=D(X)+[E(X)]^2=\lambda+\lambda^2$。由 $\lambda^2+\lambda-6=0$ 得 $\lambda=2$ 或 $\lambda=-3$，舍去负值，$\lambda=2$。

(3) 均匀分布 $E(X)=\dfrac{a+b}{2}=2$，$D(X)=\dfrac{(b-a)^2}{12}=\dfrac43$。于是 $a+b=4$，$(b-a)^2=16$，即 $b-a=4$（因 $b>a$）。解得 $a=0$，$b=4$。

(4) 指数分布 $E(X)=\dfrac1\lambda$，$E(X^2)=D(X)+[E(X)]^2=\dfrac{2}{\lambda^2}$。由 $\dfrac{2}{\lambda^2}=8$ 得 $\lambda^2=\dfrac14$，结合 $\lambda>0$ 得 $\lambda=\dfrac12$，从而 $D(X)=\dfrac{1}{\lambda^2}=4$。

\chapfourans{(1) $n=10$，$p=0.6$；(2) $\lambda=2$；(3) $a=0$，$b=4$；(4) $\lambda=\dfrac12$，$D(X)=4$}

\else \answerarea[6] \fi

\textbf{5.} 设二维离散型随机变量 $(X,Y)$ 的联合分布律为
\[
\begin{array}{c|ccc}
X\backslash Y & 0 & 1 & 2\\
\hline
0 & 0.15 & 0.05 & 0.05\\
1 & 0.05 & 0.40 & 0.05\\
2 & 0.05 & 0.05 & 0.15
\end{array}
\]
\begin{enumerate}[label=(\arabic*)]
    \item 求 $X,Y$ 的边缘分布及 $E(X),D(X),E(Y),D(Y)$；
    \item 求 $\Cov(X,Y)$ 与相关系数 $\rho_{XY}$；
    \item 判断 $X$ 与 $Y$ 是否相互独立，并写出 $(X,Y)$ 的协方差矩阵。
\end{enumerate}

\ifshowsolutions%
\boxed{\textbf{解}}
\textbf{思路：}边缘分布按行、列求和；协方差用 $\Cov(X,Y)=E(XY)-E(X)E(Y)$，其中 $E(XY)$ 要把每个格子上的 $x_iy_jp_{ij}$ 全部相加；判断独立性只需找到一个格子使 $p_{ij}\ne p_{i\cdot}p_{\cdot j}$。

(1) 各行求和得 $P(X=0)=0.25$，$P(X=1)=0.50$，$P(X=2)=0.25$；各列求和得 $P(Y=0)=0.25$，$P(Y=1)=0.50$，$P(Y=2)=0.25$。两者形状相同：
\[
E(X)=0(0.25)+1(0.50)+2(0.25)=1,
\]
\[
E(X^2)=0(0.25)+1(0.50)+4(0.25)=1.5,\qquad D(X)=1.5-1^2=0.5.
\]
同理由对称性 $E(Y)=1$，$E(Y^2)=1.5$，$D(Y)=0.5$。

(2)
\[
\begin{aligned}
E(XY)&=\sum_{i}\sum_{j}x_iy_jp_{ij}\\
&=1\cdot1\cdot0.40+1\cdot2\cdot0.05+2\cdot1\cdot0.05+2\cdot2\cdot0.15\\
&=0.40+0.10+0.10+0.60=1.2.
\end{aligned}
\]
\[
\Cov(X,Y)=E(XY)-E(X)E(Y)=1.2-1\times1=0.2.
\]
\[
\rho_{XY}=\frac{\Cov(X,Y)}{\sqrt{D(X)D(Y)}}=\frac{0.2}{\sqrt{0.5\times0.5}}=\frac{0.2}{0.5}=0.4.
\]

(3) 取 $i=0,j=0$：$p_{00}=0.15$，而 $P(X=0)P(Y=0)=0.25\times0.25=0.0625$，两者不等，故 $X$ 与 $Y$ 不独立。

协方差矩阵为
\[
\Sigma=\begin{pmatrix}D(X)&\Cov(X,Y)\\\Cov(X,Y)&D(Y)\end{pmatrix}
=\begin{pmatrix}0.5&0.2\\0.2&0.5\end{pmatrix}.
\]

\chapfourans{(1) $E(X)=E(Y)=1$，$D(X)=D(Y)=0.5$；(2) $\Cov(X,Y)=0.2$，$\rho_{XY}=0.4$；(3) 不独立，$\Sigma=\begin{pmatrix}0.5&0.2\\0.2&0.5\end{pmatrix}$}

\else \answerarea[10] \fi

\subsection{综合提高题}

\textbf{6.} 设二维随机变量 $(X,Y)$ 在区域 $D=\{(x,y):0<y<x<1\}$ 上服从均匀分布。
\begin{enumerate}[label=(\arabic*)]
    \item 写出联合密度及 $X,Y$ 的边缘密度；
    \item 求 $E(X),E(Y),D(X),D(Y)$；
    \item 求 $\Cov(X,Y)$ 与 $\rho_{XY}$，并判断 $X$ 与 $Y$ 是否独立。
\end{enumerate}

\ifshowsolutions%
\boxed{\textbf{解}}
\textbf{思路：}均匀分布的密度等于区域面积的倒数；求二维期望就是“被积函数乘联合密度”在区域上积分。区域不是矩形，边缘密度要按另一变量在该区域内的取值范围定积分。

(1) 区域 $D$ 的面积为 $\dfrac12$，故
\[
f(x,y)=\begin{cases}2,&0<y<x<1,\\0,&\text{其他}.\end{cases}
\]
固定 $x\in(0,1)$，$y$ 的范围是 $(0,x)$，故
\[
f_X(x)=\int_0^x2\,\mathrm{d}y=2x,\quad 0<x<1.
\]
固定 $y\in(0,1)$，$x$ 的范围是 $(y,1)$，故
\[
f_Y(y)=\int_y^1 2\,\mathrm{d}x=2(1-y),\quad 0<y<1.
\]

(2)
\[
E(X)=\int_0^1 x\cdot2x\,\mathrm{d}x=2\int_0^1x^2\,\mathrm{d}x=\frac23,
\]
\[
E(X^2)=\int_0^1x^2\cdot2x\,\mathrm{d}x=2\int_0^1x^3\,\mathrm{d}x=\frac12,
\qquad D(X)=\frac12-\frac49=\frac1{18}.
\]
\[
E(Y)=\int_0^1y\cdot2(1-y)\,\mathrm{d}y=2\left(\frac12-\frac13\right)=\frac13,
\]
\[
E(Y^2)=\int_0^1y^2\cdot2(1-y)\,\mathrm{d}y=2\left(\frac13-\frac14\right)=\frac16,
\qquad D(Y)=\frac16-\frac19=\frac1{18}.
\]

(3) 用联合密度计算 $E(XY)$：
\[
E(XY)=\int_0^1\!\!\int_0^x2xy\,\mathrm{d}y\,\mathrm{d}x
=\int_0^1x\cdot x^2\,\mathrm{d}x=\frac14.
\]
\[
\Cov(X,Y)=\frac14-\frac23\cdot\frac13=\frac14-\frac29=\frac1{36},
\qquad
\rho_{XY}=\frac{1/36}{\sqrt{(1/18)(1/18)}}=\frac{1/36}{1/18}=\frac12.
\]
独立性：在 $0<y<x<1$ 上 $f(x,y)=2$，而 $f_X(x)f_Y(y)=4x(1-y)\ne2$，故不独立。事实上 $Y<X$ 恒成立，已经说明二者不独立。

\textbf{评注：}$\rho_{XY}=\dfrac12$ 并不算强相关，尽管 $Y$ 总是小于 $X$。这再次说明相关系数只反映线性关系的强弱，不能用来判断一般的函数依赖。

\chapfourans{(1) $f(x,y)=2\ (0<y<x<1)$，$f_X(x)=2x$，$f_Y(y)=2(1-y)$；(2) $E(X)=\dfrac23$，$E(Y)=\dfrac13$，$D(X)=D(Y)=\dfrac1{18}$；(3) $\Cov(X,Y)=\dfrac1{36}$，$\rho_{XY}=\dfrac12$，不独立}

\else \answerarea[10] \fi

\textbf{7.} 设二维离散型随机变量 $(X,Y)$ 的联合分布律为
\[
\begin{array}{c|ccc}
X\backslash Y & 0 & 1 & 2\\
\hline
0 & 0.2 & 0.1 & 0.3\\
1 & 0.1 & 0.2 & 0.1
\end{array}
\]
\begin{enumerate}[label=(\arabic*)]
    \item 求条件期望 $E(X\mid Y=0)$、$E(X\mid Y=1)$、$E(X\mid Y=2)$；
    \item 验证重期望公式 $E(X)=E[E(X\mid Y)]$；
    \item 求 $D(X)$，并用方差分解公式 $D(X)=E[D(X\mid Y)]+D[E(X\mid Y)]$ 验证。
\end{enumerate}

\ifshowsolutions%
\boxed{\textbf{解}}
\textbf{思路：}先算 $Y$ 的边缘分布，在 $Y=y$ 的条件下把该列的联合概率除以列和得到条件分布；条件期望是 $y$ 的函数，记为 $Z=E(X\mid Y)$ 后，再对 $Y$ 的分布求期望和方差。

(1) 列和为 $P(Y=0)=0.3$，$P(Y=1)=0.3$，$P(Y=2)=0.4$。在 $Y=0$ 下，$X$ 取 $0,1$ 的条件概率分别为 $\dfrac{0.2}{0.3}=\dfrac23$、$\dfrac{0.1}{0.3}=\dfrac13$，故
\[
E(X\mid Y=0)=0\cdot\frac23+1\cdot\frac13=\frac13.
\]
同理
\[
E(X\mid Y=1)=\frac{0.1}{0.3}\cdot0+\frac{0.2}{0.3}\cdot1=\frac23,
\qquad
E(X\mid Y=2)=\frac{0.3}{0.4}\cdot0+\frac{0.1}{0.4}\cdot1=\frac14.
\]
即
\[
Z=E(X\mid Y)=\begin{cases}\frac13,&Y=0,\\[2pt]\frac23,&Y=1,\\[2pt]\frac14,&Y=2.\end{cases}
\]

(2) $X$ 的边缘分布为 $P(X=0)=0.6$，$P(X=1)=0.4$，故 $E(X)=0.4$。另一方面
\[
E(Z)=\frac13(0.3)+\frac23(0.3)+\frac14(0.4)=0.1+0.2+0.1=0.4,
\]
与 $E(X)$ 相同，重期望公式成立。

(3) 先算 $D(X)$：$E(X^2)=0^2(0.6)+1^2(0.4)=0.4$，所以
\[
D(X)=0.4-0.4^2=0.24=\frac{6}{25}.
\]
条件方差：给定 $Y=y$ 时 $X$ 是参数为 $q$ 的两点分布，故 $D(X\mid Y=y)=q(1-q)$。于是
\[
D(X\mid Y=0)=\frac13\cdot\frac23=\frac29,\quad
D(X\mid Y=1)=\frac23\cdot\frac13=\frac29,\quad
D(X\mid Y=2)=\frac14\cdot\frac34=\frac3{16}.
\]
\[
E[D(X\mid Y)]=0.3\cdot\frac29+0.3\cdot\frac29+0.4\cdot\frac3{16}
=\frac{2}{15}+\frac{3}{40}=\frac{25}{120}=\frac5{24}.
\]
再算 $D(Z)$：$E(Z^2)=0.3\cdot\dfrac19+0.3\cdot\dfrac49+0.4\cdot\dfrac1{16}
=\dfrac1{30}+\dfrac2{15}+\dfrac1{40}=\dfrac{23}{120}$，故
\[
D(Z)=\frac{23}{120}-\left(\frac25\right)^2=\frac{23}{120}-\frac{4}{25}=\frac{115-96}{600}=\frac{19}{600}.
\]
两项相加：
\[
E[D(X\mid Y)]+D(Z)=\frac5{24}+\frac{19}{600}=\frac{125+19}{600}=\frac{144}{600}=\frac6{25}=D(X),
\]
与 (3) 第一段的结果一致。

\chapfourans{(1) $E(X\mid Y=0)=\dfrac13$，$E(X\mid Y=1)=\dfrac23$，$E(X\mid Y=2)=\dfrac14$；(2) $E[E(X\mid Y)]=0.4=E(X)$；(3) $D(X)=\dfrac6{25}$，且 $\dfrac5{24}+\dfrac{19}{600}=\dfrac6{25}$}

\else \answerarea[10] \fi

\textbf{8.} 设 $E(X)=0$，$E(Y)=1$，$D(X)=4$，$D(Y)=1$，$\rho_{XY}=\dfrac12$。
\begin{enumerate}[label=(\arabic*)]
    \item 求 $\Cov(X,Y)$、$D(X+Y)$、$D(X-Y)$；
    \item 求 $\Cov(X+Y,\,X-Y)$ 与 $\rho(X+Y,\,X-Y)$；
    \item 求常数 $a,b$ 使 $E[(Y-aX-b)^2]$ 最小，并求该最小值。
\end{enumerate}

\ifshowsolutions%
\boxed{\textbf{解}}
\textbf{思路：}由 $\rho$ 的定义先反解协方差；方差的和差公式是 $D(X\pm Y)=D(X)+D(Y)\pm2\Cov(X,Y)$；第 (3) 问是均方意义下的最优线性预测，把目标函数配方即可，不必解多元微积分。

(1) 由 $\rho_{XY}=\dfrac{\Cov(X,Y)}{\sqrt{D(X)}\sqrt{D(Y)}}$ 得
\[
\Cov(X,Y)=\rho_{XY}\sqrt{D(X)}\sqrt{D(Y)}=\frac12\cdot2\cdot1=1.
\]
于是
\[
D(X+Y)=4+1+2(1)=7,\qquad D(X-Y)=4+1-2(1)=3.
\]

(2) 利用协方差的双线性性：
\[
\Cov(X+Y,X-Y)=\Cov(X,X)-\Cov(X,Y)+\Cov(Y,X)-\Cov(Y,Y)
=4-1+1-1=3.
\]
相关系数为
\[
\rho(X+Y,X-Y)=\frac{3}{\sqrt{D(X+Y)D(X-Y)}}=\frac{3}{\sqrt{7\cdot3}}=\frac{3}{\sqrt{21}}=\sqrt{\frac37}.
\]

(3) 把目标函数展开：
\[
E[(Y-aX-b)^2]=D(Y-aX)+[E(Y-aX)-b]^2.
\]
先选 $b=E(Y-aX)=1-aE(X)=1$，使第二项为零；再由
\[
D(Y-aX)=D(Y)-2a\Cov(X,Y)+a^2D(X)=1-2a+4a^2
\]
对 $a$ 配方：
\[
4a^2-2a+1=4\left(a-\frac14\right)^2+\frac34.
\]
故当 $a=\dfrac14$、$b=1$ 时平方误差最小，最小值为 $\dfrac34$。这也与公式 $D(Y)(1-\rho_{XY}^2)=1\cdot\dfrac34$ 一致。

\chapfourans{(1) $\Cov(X,Y)=1$，$D(X+Y)=7$，$D(X-Y)=3$；(2) $\Cov(X+Y,X-Y)=3$，$\rho(X+Y,X-Y)=\sqrt{\dfrac37}$；(3) $a=\dfrac14$，$b=1$，最小值为 $\dfrac34$}

\else \answerarea[10] \fi

\textbf{9.} 设连续型随机变量 $X$ 的概率密度为
\[
f(x)=A\mathrm{e}^{-x^2+2x},\quad -\infty<x<+\infty.
\]
\begin{enumerate}[label=(\arabic*)]
    \item 求常数 $A$，并指出 $X$ 服从什么分布；
    \item 求 $E(X)$、$D(X)$ 与 $E(X^2)$；
    \item 求三阶中心矩 $E[(X-E(X))^3]$ 与四阶中心矩 $E[(X-E(X))^4]$。
\end{enumerate}

\ifshowsolutions%
\boxed{\textbf{解}}
\textbf{思路：}指数上配方，把密度化成标准正态密度的形式，再与 $N(\mu,\sigma^2)$ 的密度对照读出 $\mu,\sigma^2$。正态分布关于均值对称，故奇数阶中心矩为零；四阶中心矩用标准化变量 $Z\sim N(0,1)$ 的 $E(Z^4)=3$。

(1) 配方：
\[
-x^2+2x=-(x-1)^2+1,
\]
故
\[
f(x)=A\mathrm{e}\cdot \mathrm{e}^{-(x-1)^2}.
\]
而 $N(\mu,\sigma^2)$ 的密度为 $\dfrac{1}{\sqrt{2\pi}\sigma}\exp\left(-\dfrac{(x-\mu)^2}{2\sigma^2}\right)$。对照得 $\mu=1$，$\dfrac{1}{2\sigma^2}=1$ 即 $\sigma^2=\dfrac12$。于是
\[
A\mathrm{e}=\frac{1}{\sqrt{2\pi}\sqrt{1/2}}=\frac{1}{\sqrt{\pi}},
\qquad A=\frac{1}{\mathrm{e}\sqrt{\pi}}.
\]
即 $X\sim N\!\left(1,\dfrac12\right)$。

(2) 由正态分布的数字特征直接得
\[
E(X)=\mu=1,\qquad D(X)=\sigma^2=\frac12,\qquad E(X^2)=D(X)+[E(X)]^2=\frac12+1=\frac32.
\]

(3) 令 $Z=\sqrt2\,(X-1)\sim N(0,1)$，则 $X-1=\dfrac{Z}{\sqrt2}$。正态分布关于均值对称，故
\[
E[(X-1)^3]=\frac{1}{2\sqrt2}E(Z^3)=0.
\]
又 $E(Z^4)=3$（可由 $E(Z^2)=1$ 与正态矩公式 $E(Z^{2k})=(2k-1)!!$ 得到），所以
\[
E[(X-1)^4]=\frac{1}{(\sqrt2)^4}E(Z^4)=\frac{3}{4}.
\]

\chapfourans{(1) $A=\dfrac{1}{\mathrm{e}\sqrt{\pi}}$，$X\sim N\!\left(1,\dfrac12\right)$；(2) $E(X)=1$，$D(X)=\dfrac12$，$E(X^2)=\dfrac32$；(3) 三阶中心矩为 $0$，四阶中心矩为 $\dfrac34$}

\else \answerarea[10] \fi

\subsection{真题风格与综合训练}

\textbf{10.} 某投资者把全部资金按比例 $w$ 投入资产 $X$、按比例 $1-w$ 投入资产 $Y$（$0\le w\le1$），组合收益率为 $R=wX+(1-w)Y$。已知 $E(X)=0.12$，$E(Y)=0.06$，$D(X)=D(Y)=\dfrac1{25}$，$\Cov(X,Y)=\dfrac1{100}$。
\begin{enumerate}[label=(\arabic*)]
    \item 求 $E(R)$ 关于 $w$ 的表达式；
    \item 求 $D(R)$ 关于 $w$ 的表达式，并求使 $D(R)$ 最小的 $w$ 及最小方差；
    \item 若要求 $E(R)=0.10$，求相应的 $w$，并求此时的 $D(R)$，与 (2) 的最小方差比较。
\end{enumerate}

\ifshowsolutions%
\boxed{\textbf{解}}
\textbf{思路：}期望用线性性直接代入；方差是 $w$ 的二次函数，展开后配方或求导找顶点。注意“收益目标”约束的是期望，与“风险最小”是两个不同的问题。

(1) 由线性性
\[
E(R)=wE(X)+(1-w)E(Y)=0.12w+0.06(1-w)=0.06+0.06w.
\]

(2) 由方差的和差公式
\[
D(R)=w^2D(X)+(1-w)^2D(Y)+2w(1-w)\Cov(X,Y).
\]
代入 $D(X)=D(Y)=\dfrac1{25}$、$\Cov(X,Y)=\dfrac1{100}$：
\[
D(R)=\frac1{25}\bigl[w^2+(1-w)^2\bigr]+\frac1{50}w(1-w).
\]
利用 $w^2+(1-w)^2=1-2w+2w^2$ 展开：
\[
D(R)=\frac1{25}(1-2w+2w^2)+\frac1{50}(w-w^2)
=\frac1{25}-\frac{3}{50}w+\frac{3}{50}w^2.
\]
这是开口向上的二次函数，对称轴为
\[
w=\frac{3/50}{2\cdot3/50}=\frac12.
\]
故 $w=\dfrac12$ 时方差最小，最小值为
\[
D(R)\Big|_{w=1/2}=\frac1{25}-\frac{3}{100}+\frac{3}{200}
=\frac{8-6+3}{200}=\frac{5}{200}=\frac1{40}.
\]
此时 $E(R)=0.09$。

(3) 由 $0.06+0.06w=0.10$ 得 $w=\dfrac{0.04}{0.06}=\dfrac23$。代入方差表达式：
\[
D(R)\Big|_{w=2/3}
=\frac1{25}-\frac{3}{50}\cdot\frac23+\frac{3}{50}\cdot\frac49
=\frac1{25}-\frac1{25}+\frac{2}{75}
=\frac{2}{75}.
\]
因为 $\dfrac{2}{75}\approx0.0267>\dfrac1{40}=0.025$，为了把期望收益从 $0.09$ 提高到 $0.10$，必须承担更大的方差，这正是收益与风险的权衡。

\chapfourans{(1) $E(R)=0.06+0.06w$；(2) $D(R)=\dfrac1{25}-\dfrac{3}{50}w+\dfrac{3}{50}w^2$，$w=\dfrac12$ 时最小，最小方差为 $\dfrac1{40}$；(3) $w=\dfrac23$，$D(R)=\dfrac{2}{75}>\dfrac1{40}$}

\else \answerarea[10] \fi

\textbf{11.} 设 $(X,Y)$ 服从二维正态分布，$E(X)=E(Y)=0$，$D(X)=D(Y)=1$，$\rho_{XY}=\dfrac12$。
\begin{enumerate}[label=(\arabic*)]
    \item 求 $D(X+Y)$、$D(X-Y)$ 及 $\Cov(X+Y,X-Y)$；
    \item 判断 $X+Y$ 与 $X-Y$ 是否相互独立，说明理由；
    \item 求 $E[(X+Y)^2(X-Y)^2]$ 与 $P(|X+Y|\le\sqrt3)$（已知 $\Phi(1)=0.8413$）。
\end{enumerate}

\ifshowsolutions%
\boxed{\textbf{解}}
\textbf{思路：}正态变量的线性组合仍服从正态分布，且联合正态下“不相关即独立”。第 (3) 问先利用独立性把乘积的期望拆开，再把 $X+Y$ 标准化后查标准正态分布表。

(1) $\Cov(X,Y)=\rho_{XY}\sqrt{D(X)}\sqrt{D(Y)}=\dfrac12$，故
\[
D(X+Y)=1+1+2\cdot\frac12=3,\qquad
D(X-Y)=1+1-2\cdot\frac12=1.
\]
\[
\Cov(X+Y,X-Y)=D(X)-D(Y)=1-1=0.
\]

(2) 相互独立。因为 $(X,Y)$ 服从二维正态分布，$(X+Y,\,X-Y)$ 是 $(X,Y)$ 的线性变换，仍服从二维正态分布；由 (1) 知二者协方差为零，即不相关。对二维正态变量，不相关与独立等价，所以 $X+Y$ 与 $X-Y$ 相互独立。

(3) 由独立性，
\[
E[(X+Y)^2(X-Y)^2]=E[(X+Y)^2]\cdot E[(X-Y)^2].
\]
而 $E(X+Y)=0$，$E(X-Y)=0$，故
\[
E[(X+Y)^2]=D(X+Y)=3,\qquad E[(X-Y)^2]=D(X-Y)=1,
\]
从而乘积的期望为 $3$。

又 $X+Y\sim N(0,3)$，标准化得 $Z=\dfrac{X+Y}{\sqrt3}\sim N(0,1)$，于是
\[
P(|X+Y|\le\sqrt3)=P(|Z|\le1)=2\Phi(1)-1=2(0.8413)-1=0.6826.
\]

\chapfourans{(1) $D(X+Y)=3$，$D(X-Y)=1$，$\Cov(X+Y,X-Y)=0$；(2) 独立；(3) $E[(X+Y)^2(X-Y)^2]=3$，$P(|X+Y|\le\sqrt3)\approx0.6826$}

\else \answerarea[10] \fi

\textbf{12.} 设 $X\sim U(0,1)$，在给定 $X=x$（$0<x<1$）的条件下，$Y$ 服从区间 $(0,x)$ 上的均匀分布。
\begin{enumerate}[label=(\arabic*)]
    \item 求 $(X,Y)$ 的联合密度与 $Y$ 的边缘密度；
    \item 求 $E(Y)$、$D(Y)$；
    \item 求 $\Cov(X,Y)$ 与 $\rho_{XY}$；
    \item 求条件期望 $E(X\mid Y=y)$，$0<y<1$。
\end{enumerate}

\ifshowsolutions%
\boxed{\textbf{解}}
\textbf{思路：}$f(x,y)=f_X(x)f_{Y\mid X}(y\mid x)$，注意条件密度只在 $0<y<x$ 上非零。求 $Y$ 的数字特征既可以直接用边缘密度积分，也可以先算条件矩再用重期望公式；第 (4) 问先求 $f_Y$，再由 $f_{X\mid Y}=f/f_Y$ 计算。

(1) $f_X(x)=1$（$0<x<1$），$f_{Y\mid X}(y\mid x)=\dfrac1x$（$0<y<x$），故
\[
f(x,y)=\begin{cases}\dfrac1x,&0<y<x<1,\\[4pt]0,&\text{其他}.\end{cases}
\]
固定 $y\in(0,1)$，$x$ 的范围为 $(y,1)$，于是
\[
f_Y(y)=\int_y^1\frac1x\,\mathrm{d}x=-\ln y,\quad 0<y<1.
\]

(2) 先求条件矩：$E(Y\mid X=x)=\dfrac{x}{2}$，$D(Y\mid X=x)=\dfrac{x^2}{12}$。由重期望公式
\[
E(Y)=E\!\left(\frac{X}{2}\right)=\frac12\cdot\frac12=\frac14.
\]
再求 $E(Y^2)$：
\[
E(Y^2)=E[E(Y^2\mid X)]
=E\!\left[\frac{X^2}{3}\right]=\frac13\cdot\frac13=\frac19.
\]
故
\[
D(Y)=\frac19-\frac1{16}=\frac{16-9}{144}=\frac7{144}.
\]

(3)
\[
E(XY)=E[E(XY\mid X)]=E[X\,E(Y\mid X)]
=E\!\left[\frac{X^2}{2}\right]=\frac16.
\]
又 $E(X)=\dfrac12$，$E(Y)=\dfrac14$，故
\[
\Cov(X,Y)=\frac16-\frac12\cdot\frac14=\frac16-\frac18=\frac1{24}.
\]
$X\sim U(0,1)$ 给出 $D(X)=\dfrac1{12}$，于是
\[
\rho_{XY}=\frac{1/24}{\sqrt{\frac1{12}\cdot\frac7{144}}}
=\frac{1/24}{\sqrt{\frac7{1728}}}
=\frac{1}{24}\cdot\frac{24\sqrt3}{\sqrt7}
=\sqrt{\frac37}\approx0.655.
\]

(4) 当 $0<y<x<1$ 时，
\[
f_{X\mid Y}(x\mid y)=\frac{f(x,y)}{f_Y(y)}=\frac{1}{x(-\ln y)}.
\]
因此
\[
E(X\mid Y=y)=\int_y^1 x\cdot\frac{1}{x(-\ln y)}\,\mathrm{d}x
=\frac{1}{-\ln y}\int_y^1\mathrm{d}x
=\frac{1-y}{-\ln y}.
\]
可以检验：$\displaystyle\int_0^1\frac{1-y}{-\ln y}(-\ln y)\,\mathrm{d}y=\int_0^1(1-y)\,\mathrm{d}y=\frac12=E(X)$，与重期望公式相符。

\chapfourans{(1) $f(x,y)=\dfrac1x\ (0<y<x<1)$，$f_Y(y)=-\ln y\ (0<y<1)$；(2) $E(Y)=\dfrac14$，$D(Y)=\dfrac7{144}$；(3) $\Cov(X,Y)=\dfrac1{24}$，$\rho_{XY}=\sqrt{\dfrac37}$；(4) $E(X\mid Y=y)=\dfrac{1-y}{-\ln y}$}

\else \answerarea[10] \fi
